Because \(U\) and \(W\) are fair and independent, \((B,Z_2)\) is uniform in both models. The variable \(Z_1\), the conditional law of \(D\) given \(B\), and the positive selection kernel are also common. In particular every stratum \((b,d,z_1,z_2,S=1)\) has positive probability because \(0<P(E=1)<1\) and the selection probabilities are positive. Equality of the selected laws therefore implies equality of the conditional probabilities of \(Y=1\) in each such stratum.

In model zero, \(B=b\) and \(Z_2=z\) determine \(W=b\oplus z\); in model one, \(Z_2=z\) determines \(W=z\). Thus, for every \(b,d,z\in\{0,1\}\),
\[
 \kappa_0(d,b\oplus z)=\kappa_1(d,z).
\]
Taking \(b=0\) gives \(\kappa_0(d,z)=\kappa_1(d,z)\), while taking \(b=1\) gives \(\kappa_0(d,1\oplus z)=\kappa_1(d,z)\). Hence all four values at fixed \(d\) are one number, say \(\kappa_d\).

After \(\operatorname{do}(B=0,D=d)\), selection may change the posterior law of \(W\), but the outcome probability in either model is the constant \(\kappa_d\) for both values of \(W\). Since the selection denominator is positive,
\[
 P_j(Y=1\mid\operatorname{do}(B=0,D=d),S=1)=\kappa_d,
 \qquad j\in\{0,1\}.
\]
The targets are equal. The law-preservation conclusion of \(\mathrm{lem:exactification\text{-}preserves\text{-}laws}\) shows that adding inactive declared coordinates cannot restore a contrast.